% GENERATED FILE. Edit canonical lessons, then run npm run generate:books.
% canonical-source-hash: fnv1a64:a6ea62f37f9f9a82
% canonical-lessons: FA-C35-khordan, FA-C35-nushidan, FA-C35-didan, FA-C35-shenidan, FA-C35-neshastan

\chapter{To eat, To drink, To see, To hear, To sit}
\label{ch:fa-to-eat-to-drink-to-see-to-hear-to-sit}
\hlchaptermodality{35}

\begin{chapteropening}
\textbf{By the end of this chapter:} \emph{I can say to eat, to drink, to see, to hear and to sit.}

Complete the last lesson of chapter 35: I can say to eat, to drink, to see, to hear and to sit.
\end{chapteropening}

\section[\texorpdfstring{\fa{خوردن} (khordan)}{خوردن (khordan)}]{\fa{خوردن} (khordan) — to eat}

\label{lesson:FA-C35-khordan}



\begin{quote}
Before the new one: say the Persian for a gentleman; Mr, then the Persian for a lady; Mrs.
\end{quote}

\subsection*{You'll want to know: \fa{خوردن}}
\textbf{\fa{خوردن}} — \emph{khordan} — \textquotedblleft{}to eat\textquotedblright{}.

\begin{cousinweb}
\emph{khordan} is a native Persian verb, inherited rather than borrowed: EtymDB traces it to the Middle Persian stem \emph{xwar-}, \textquotedblleft{}to eat\textquotedblright{}, and further to Proto-Iranian \emph{*hwar-}, \textquotedblleft{}to eat, to consume\textquotedblright{}.
\end{cousinweb}

\begin{grammarlens}[title={What you've built}]
One more word for this chapter.
\end{grammarlens}

\subsection*{Your turn}
\begin{itemize}
\raggedright
  \item \emph{Say it:} \emph{khordan}
  \item \emph{Say it:} \emph{khordan}, once more
  \item \emph{Say it:} say \emph{khordan}
  \item \emph{Recall:} say the Persian for a gentleman; Mr, then the Persian for a lady; Mrs, then say \emph{khordan} again
\end{itemize}

\begin{tcolorbox}[breakable,colback=teal!4,colframe=teal!35!black,title={Before you move on}]
What does \fa{خوردن} mean? (\textquotedblleft{}To eat\textquotedblright{}.) Say it once more.
\end{tcolorbox}
\section[\texorpdfstring{\fa{نوشیدن} (nushidan)}{نوشیدن (nushidan)}]{\fa{نوشیدن} (nushidan) — to drink}

\label{lesson:FA-C35-nushidan}



\begin{quote}
Before the new one: say the Persian for a lady; Mrs, then the Persian for to eat.
\end{quote}

\subsection*{You'll want to know: \fa{نوشیدن}}
\textbf{\fa{نوشیدن}} — \emph{nushidan} — \textquotedblleft{}to drink\textquotedblright{}.

\begin{cousinweb}
\emph{nushidan} is a Persian verb built on \emph{nuš}, \textquotedblleft{}sweet drink\textquotedblright{}, with the verb ending \emph{-idan}. \emph{nuš} is usually traced to Middle Persian \emph{anōš}, \textquotedblleft{}immortality, ambrosia\textquotedblright{}.
\end{cousinweb}

\begin{grammarlens}[title={What you've built}]
One more word for this chapter.
\end{grammarlens}

\subsection*{Your turn}
\begin{itemize}
\raggedright
  \item \emph{Say it:} \emph{nushidan}
  \item \emph{Say it:} \emph{nushidan}, once more
  \item \emph{Say it:} say \emph{nushidan}
  \item \emph{Recall:} say the Persian for a lady; Mrs, then the Persian for to eat, then say \emph{nushidan} again
\end{itemize}

\begin{tcolorbox}[breakable,colback=teal!4,colframe=teal!35!black,title={Before you move on}]
What does \fa{نوشیدن} mean? (\textquotedblleft{}To drink\textquotedblright{}.) Say it once more.
\end{tcolorbox}
\section[\texorpdfstring{\fa{دیدن} (didan)}{دیدن (didan)}]{\fa{دیدن} (didan) — to see}

\label{lesson:FA-C35-didan}



\begin{quote}
Before the new one: say the Persian for to eat, then the Persian for to drink.
\end{quote}

\subsection*{You'll want to know: \fa{دیدن}}
\textbf{\fa{دیدن}} — \emph{didan} — \textquotedblleft{}to see\textquotedblright{}.

\begin{cousinweb}
\emph{didan} is a native Persian word, inherited rather than borrowed; it is usually traced to Middle Persian \emph{dīdan}, \textquotedblleft{}to see\textquotedblright{}.
\end{cousinweb}

\begin{grammarlens}[title={What you've built}]
One more word for this chapter.
\end{grammarlens}

\subsection*{Your turn}
\begin{itemize}
\raggedright
  \item \emph{Say it:} \emph{didan}
  \item \emph{Say it:} \emph{didan}, once more
  \item \emph{Say it:} say \emph{didan}
  \item \emph{Recall:} say the Persian for to eat, then the Persian for to drink, then say \emph{didan} again
\end{itemize}

\begin{tcolorbox}[breakable,colback=teal!4,colframe=teal!35!black,title={Before you move on}]
What does \fa{دیدن} mean? (\textquotedblleft{}To see\textquotedblright{}.) Say it once more.
\end{tcolorbox}
\section[\texorpdfstring{\fa{شنیدن} (shenidan)}{شنیدن (shenidan)}]{\fa{شنیدن} (shenidan) — to hear}

\label{lesson:FA-C35-shenidan}



\begin{quote}
Before the new one: say the Persian for to drink, then the Persian for to see.
\end{quote}

\subsection*{You'll want to know: \fa{شنیدن}}
\textbf{\fa{شنیدن}} — \emph{shenidan} — \textquotedblleft{}to hear\textquotedblright{}.

\begin{cousinweb}
\emph{shenidan} is a native Persian word, inherited rather than borrowed; it is usually traced to Middle Persian \emph{ašnūdan}, \textquotedblleft{}to hear\textquotedblright{}.
\end{cousinweb}

\begin{grammarlens}[title={What you've built}]
One more word for this chapter.
\end{grammarlens}

\subsection*{Your turn}
\begin{itemize}
\raggedright
  \item \emph{Say it:} \emph{shenidan}
  \item \emph{Say it:} \emph{shenidan}, once more
  \item \emph{Say it:} say \emph{shenidan}
  \item \emph{Recall:} say the Persian for to drink, then the Persian for to see, then say \emph{shenidan} again
\end{itemize}

\begin{tcolorbox}[breakable,colback=teal!4,colframe=teal!35!black,title={Before you move on}]
What does \fa{شنیدن} mean? (\textquotedblleft{}To hear\textquotedblright{}.) Say it once more.
\end{tcolorbox}
\section[\texorpdfstring{\fa{نشستن} (neshastan)}{نشستن (neshastan)}]{\fa{نشستن} (neshastan) — to sit}

\label{lesson:FA-C35-neshastan}



\begin{quote}
Before the new one: say the Persian for to see, then the Persian for to hear.
\end{quote}

\subsection*{You'll want to know: \fa{نشستن}}
\textbf{\fa{نشستن}} — \emph{neshastan} — \textquotedblleft{}to sit\textquotedblright{}.

\begin{cousinweb}
\emph{neshastan} is a native Persian word, inherited rather than borrowed; it is usually traced to Middle Persian \emph{nišastan}, \textquotedblleft{}to sit\textquotedblright{}. It is usually analysed as \emph{ni-}, \textquotedblleft{}down\textquotedblright{}, on the Indo-European root that also gives English \emph{sit}.
\end{cousinweb}

\begin{grammarlens}[title={What you've built}]
One more word for this chapter.
\end{grammarlens}

\subsection*{Your turn}
\begin{itemize}
\raggedright
  \item \emph{Say it:} \emph{neshastan}
  \item \emph{Say it:} \emph{neshastan}, once more
  \item \emph{Say it:} say \emph{neshastan}
  \item \emph{Recall:} say the Persian for to see, then the Persian for to hear, then say \emph{neshastan} again
\end{itemize}

\begin{tcolorbox}[breakable,colback=teal!4,colframe=teal!35!black,title={Before you move on}]
What does \fa{نشستن} mean? (\textquotedblleft{}To sit\textquotedblright{}.) Say it once more.
\end{tcolorbox}
